\section*{序言（食用指南）}


本讲义主要参考北大版《高等数学》（第三版），
按照高等数学B课程要求编写，
可以配合教材作为学习和复习的辅助材料使用.
讲义分为十二章，
与教材的章节划分一致.
每一章主要分为“\textbf{复习要点}”和“\textbf{解题方法}”两大部分.
\vspace{1em}

“\textbf{复习要点}”主要安排以下内容：

(1)~~构建知识框架；

(2)~~梳理需要集中背诵的定理、公式等；

(3)~~深入展开本章需要加强理解的知识点、辨析易混淆的概念、提醒易错点；

(4)~~拓展相关知识（小节标题以“$*$”标注，若不感兴趣可自由跳过）.\vspace{1em}

“\textbf{解题方法}”主要安排以下内容：

(1)~~重点题型的解题方法梳理
（题型是否“重点”是以北京大学高数B期中/期末考试为标准的）；

(2)~~例题和解答过程，有时附有分析，展示思考过程；

(3)~~进阶的解题方法和例题（小节标题以“*”标注）\vspace{1em}

这份讲义并不像教科书一样从零开始构建高等数学的知识体系，因此只能用作学习的辅助材
料. 如果您的目标是通过学校的高数期中/期末考试，或者纯粹地希望学好高等数学，那么可以使
用这份讲义；如果您的目标是考研或竞赛，还请使用专用的相关教辅.

这份讲义的编写和使用历经了一年的时间. 在此期间，我的数学水平、表达能力和排版技术
还处于进步的过程中，也不断发现此前已经落笔的章节还有许多错漏和缺憾. 在各路读者的帮助
下，我会在后续的版本中继续修正错误. 另外，讲义设计的初衷是帮助北大生科院的大一同学们
无痛学习高数，因此编写的细节处可能没有为更广大读者的方便着想. 诚恳地希望各路读者对讲
义提出宝贵意见，以便改进.

这份讲义能够编写完成，离不开许多人的帮助和共同努力. 考虑到个人隐私，这里不一一列
举和致谢了. 在这里，我想重点感谢看过我讲义半成品的网友们. 他们之中，有已经学过高数的，
对讲义的内容提出了许多建议；也有正在学习高数的，冒着高数成绩被我耽误的风险充当“小白
鼠”，并向我反馈学习体验. 衷心感谢各位朋友的帮助和鼓励！


\begin{flushright}
    2025.8.3
\end{flushright}